\documentclass[handout]{beamer}
\mode<presentation>
\usetheme[secheader]{Boadilla}
\usecolortheme{whale}

\usepackage[utf8]{inputenc}
\usepackage[T1]{fontenc}
\usepackage[indonesian]{babel}
\usepackage{lmodern}
\usepackage{graphicx}
\usepackage{bookmark}

\setbeamertemplate{footline}
{
	\leavevmode%
	\hbox{%
		\begin{beamercolorbox}[wd=.333333\paperwidth,ht=2.25ex,dp=1ex,center]{author in head/foot}%
			\usebeamerfont{author in head/foot}\insertshortauthor%
		\end{beamercolorbox}%
		\begin{beamercolorbox}[wd=.433333\paperwidth,ht=2.25ex,dp=1ex,center]{title in head/foot}%
			\usebeamerfont{title in head/foot}\insertshorttitle
		\end{beamercolorbox}%
		\begin{beamercolorbox}[wd=.233333\paperwidth,ht=2.25ex,dp=1ex,right]{date in head/foot}%
			\usebeamerfont{date in head/foot}\insertshortdate{}\hspace*{2em} \insertframenumber{} / \inserttotalframenumber\hspace*{2ex}
	\end{beamercolorbox}}%
	\vskip0pt%
}

\title[Bab 8 -- Flexbox Grid]{Pemrograman Web}
\subtitle{Bab 8: Flexbox, Grid, dan Desain Responsif}
\author{Cecep Suwanda, S.Si., M.Kom.}
\institute{Universitas Bale Bandung}
\date{\today}

\begin{document}

% Slide Judul
\begin{frame}
	\titlepage
\end{frame}

% Outline dan CPMK
\begin{frame}{Outline Bab 8}
	\begin{block}{Sub-CPMK 2.2}
		Menerapkan CSS3 untuk styling, layout (Flexbox/Grid), dan desain responsif.
	\end{block}
	\begin{columns}[T]
		\begin{column}{0.5\textwidth}
			\begin{itemize}
				\item Normal Flow dan Positioning
				\item Flexbox: Konsep \& Praktik
				\item CSS Grid: Konsep \& Praktik
				\item Float dan Clear
				\item Overflow
				\item Media Queries
				\item Viewport dan Unit Responsif
				\item Mobile-First dan Breakpoint
				\item Styling Form
				\item Navigasi dan Menu
			\end{itemize}
		\end{column}
		\begin{column}{0.5\textwidth}
			\begin{itemize}
				\item Transisi
				\item Animasi dan Keyframes
				\item Transform
				\item CSS Custom Properties
				\item Spesifisitas dan Cascade
				\item Pseudo-class/element
				\item Aksesibilitas CSS
				\item BEM
				\item Best Practices Performa
				\item Referensi Standar
			\end{itemize}
		\end{column}
	\end{columns}
\end{frame}

% ========== SECTION 1: Normal Flow dan Positioning ==========
\begin{frame}{Normal Flow dan Positioning}
	\begin{itemize}
		\item \textbf{Normal flow}: elemen block vertikal, inline horizontal.
		\item Lima nilai \texttt{position}: \texttt{static}, \texttt{relative}, \texttt{absolute}, \texttt{fixed}, \texttt{sticky}.
	\end{itemize}
	\begin{block}{Karakteristik}
		\texttt{static} = default; \texttt{relative} = konteks untuk absolute; \texttt{absolute} = keluar flow (tooltip, dropdown); \texttt{fixed} = relatif viewport; \texttt{sticky} = hybrid.
	\end{block}
	\textbf{z-index} dan \textbf{stacking context}: hanya berlaku pada position non-static; nilai lebih tinggi di atas.
\end{frame}

% ========== SECTION 2-3: Flexbox ==========
\begin{frame}{Flexbox: Konsep Container dan Item}
	\begin{itemize}
		\item Layout \textbf{satu dimensi} (baris ATAU kolom).
		\item \texttt{display: flex} pada kontainer; anak menjadi \textbf{flex items}.
		\item \textbf{Main axis} (horizontal default) dan \textbf{cross axis}.
	\end{itemize}
	\begin{block}{Properti container}
		\texttt{flex-direction}, \texttt{justify-content}, \texttt{align-items}, \texttt{flex-wrap}, \texttt{gap}.
	\end{block}
	\textbf{Flex items}: \texttt{flex-grow}, \texttt{flex-shrink}, \texttt{flex-basis}; shorthand \texttt{flex: 1} = \texttt{1 1 0\%}.
\end{frame}

\begin{frame}{Flexbox: Praktik Layout}
	\begin{itemize}
		\item \textbf{Navbar responsif}: \texttt{justify-content: space-between}; mobile: \texttt{flex-direction: column}.
		\item \textbf{Pattern}: Centering (\texttt{justify-content + align-items: center}); Equal columns (\texttt{flex: 1}); Sticky footer (body flex column, main \texttt{flex: 1}).
	\end{itemize}
	\begin{block}{Flexbox vs Grid}
		Flexbox 1D (navbar, toolbar); Grid 2D (layout halaman). Keduanya saling melengkapi.
	\end{block}
\end{frame}

% ========== SECTION 4-5: CSS Grid ==========
\begin{frame}{CSS Grid: Konsep}
	\begin{itemize}
		\item Layout \textbf{dua dimensi} (baris DAN kolom).
		\item \texttt{display: grid}; \texttt{grid-template-columns}, \texttt{grid-template-rows}, \texttt{grid-template-areas}.
		\item Satuan \texttt{fr} (fraction); \texttt{repeat(n, ukuran)}; \texttt{minmax(min, max)}.
	\end{itemize}
	\begin{block}{Grid responsif tanpa media query}
		\texttt{grid-template-columns: repeat(auto-fill, minmax(250px, 1fr));} --- kolom menyesuaikan otomatis.
	\end{block}
	\textbf{Grid areas}: \texttt{grid-area: nama}; \texttt{grid-template-areas: "header header" "side main";}
\end{frame}

\begin{frame}{CSS Grid: Praktik Layout}
	\begin{itemize}
		\item Layout halaman: \texttt{grid-template-areas} dengan header, sidebar, main, footer.
		\item Kartu produk: \texttt{repeat(auto-fill, minmax(280px, 1fr))}.
		\item \texttt{auto-fill} vs \texttt{auto-fit}: auto-fill pertahankan track kosong; auto-fit melebur track kosong.
	\end{itemize}
\end{frame}

% ========== SECTION 6-7: Float dan Overflow ==========
\begin{frame}{Float, Clear, dan Overflow}
	\begin{block}{Float}
		\texttt{float: left/right}; \texttt{clear: both}; \texttt{display: flow-root} untuk BFC. Sudah usang untuk layout --- gunakan Flexbox/Grid.
	\end{block}
	\begin{block}{Overflow}
		\texttt{visible}, \texttt{hidden}, \texttt{scroll}, \texttt{auto}, \texttt{clip}. \texttt{overflow-x/y} terpisah. \texttt{text-overflow: ellipsis} dengan \texttt{white-space: nowrap}. Overflow $\neq$ visible menciptakan BFC.
	\end{block}
\end{frame}

% ========== SECTION 8-10: Media Queries dan Responsif ==========
\begin{frame}{Media Queries}
	\begin{itemize}
		\item \texttt{@media (kondisi) \{ CSS \}}; kondisi: \texttt{min-width}, \texttt{max-width}.
		\item Breakpoint umum: 480px, 768px, 1024px, 1200px.
		\item \texttt{prefers-color-scheme: dark}; \texttt{prefers-reduced-motion: reduce} (aksesibilitas).
		\item Breakpoint berdasarkan konten, bukan perangkat.
	\end{itemize}
\end{frame}

\begin{frame}{Viewport dan Mobile-First}
	\begin{itemize}
		\item Meta viewport: \texttt{<meta name="viewport" content="width=device-width, initial-scale=1.0">}.
		\item Unit: \texttt{vw}, \texttt{vh}, \texttt{dvh}, \texttt{svh} (mobile address bar).
		\item \textbf{Fluid typography}: \texttt{font-size: clamp(1rem, 2.5vw, 2rem);}
	\end{itemize}
	\begin{block}{Mobile-First}
		CSS base untuk mobile; tambah \texttt{@media (min-width: ...)} untuk layar besar. Lebih ringan; direkomendasikan.
	\end{block}
\end{frame}

% ========== SECTION 11-12: Form dan Navigasi ==========
\begin{frame}{Styling Form dan Navigasi}
	\begin{itemize}
		\item Form: \texttt{appearance: none}; \texttt{font-family: inherit} (form tidak mewarisi font!); indikator fokus wajib.
		\item Dropdown: \texttt{position: absolute} pada submenu; \texttt{:hover} dan \texttt{:focus-within} untuk aksesibilitas keyboard.
		\item Jangan hanya \texttt{:hover} --- pengguna keyboard tidak bisa hover.
	\end{itemize}
\end{frame}

% ========== SECTION 13-15: Transisi, Animasi, Transform ==========
\begin{frame}{Transisi, Animasi, dan Transform}
	\begin{block}{Transition}
		Perubahan halus dipicu state; \texttt{transition: property duration timing delay}. Properti interpolable saja.
	\end{block}
	\begin{block}{Animation}
		\texttt{@keyframes}; berjalan otomatis/berulang; cocok efek dekoratif. Transition vs Animation: pemicu vs otomatis.
	\end{block}
	\textbf{Transform}: \texttt{translate}, \texttt{scale}, \texttt{rotate}, \texttt{skew}; tidak mempengaruhi layout; centering: \texttt{transform: translate(-50\%, -50\%);}
\end{frame}

\begin{frame}{Performa Animasi}
	\begin{block}{GPU-friendly}
		Animasikan hanya \texttt{transform} dan \texttt{opacity} --- dioptimalkan GPU. Hindari \texttt{width}, \texttt{height}, \texttt{margin} (reflow mahal).
	\end{block}
	\texttt{prefers-reduced-motion: reduce} --- kurangi/matikan animasi untuk pengguna sensitif.
\end{frame}

% ========== SECTION 16-18: Custom Properties, Cascade, Pseudo ==========
\begin{frame}{CSS Custom Properties}
	\begin{itemize}
		\item \texttt{--nama: nilai;} di \texttt{:root}; akses \texttt{var(--nama)}.
		\item Runtime (ubah via JS/media query); cascade; fallback \texttt{var(--warna, blue)}.
		\item Design system: warna, spacing, radius di \texttt{:root}. Dark mode: override di \texttt{[data-theme="dark"]} atau \texttt{@media (prefers-color-scheme: dark)}.
	\end{itemize}
\end{frame}

\begin{frame}{Spesifisitas dan Pseudo}
	\begin{itemize}
		\item Specificity: inline $>$ ID $>$ class $>$ element; hindari \texttt{!important}. \texttt{@layer base, components, utilities} untuk prioritas eksplisit.
		\item Pseudo-class: \texttt{:nth-child}, \texttt{:has} (parent selector), \texttt{:is} (matches list).
		\item Pseudo-element: \texttt{::before}, \texttt{::after} (perlu \texttt{content}); tidak terbaca screen reader.
	\end{itemize}
\end{frame}

% ========== SECTION 19-21: Aksesibilitas, BEM, Performa ==========
\begin{frame}{Aksesibilitas, BEM, dan Performa}
	\begin{block}{Aksesibilitas CSS}
		Indikator fokus (\texttt{:focus-visible}); \texttt{prefers-reduced-motion}; kontras $\geq$ 4.5:1; \texttt{.sr-only} untuk skip link. Jangan \texttt{outline: none} tanpa pengganti!
	\end{block}
	\begin{block}{BEM}
		Block (\texttt{.card}), Element (\texttt{.card\_\_title}), Modifier (\texttt{.card--featured}).
	\end{block}
	Best practices: GPU-friendly; hindari \texttt{@import}; critical CSS; minifikasi; \texttt{will-change} dengan bijak.
\end{frame}

% ========== SECTION 22: Referensi ==========
\begin{frame}{Referensi Standar}
	\begin{itemize}
		\item \textbf{W3C}: spesifikasi resmi CSS.
		\item \textbf{MDN}: dokumentasi praktis, kompatibilitas.
		\item \textbf{web.dev}: tutorial, best practices.
		\item \textbf{CSS-Tricks}: panduan visual, tips layout.
		\item \textbf{Can I Use}: kompatibilitas browser.
	\end{itemize}
\end{frame}

% ========== Rangkuman dan Penutup ==========
\begin{frame}{Rangkuman Bab 8}
	\begin{itemize}
		\item Flexbox (1D) dan Grid (2D) fondasi layout modern; keduanya saling melengkapi.
		\item Mobile-first (\texttt{min-width}) menghasilkan CSS lebih ringan.
		\item Animasikan hanya \texttt{transform}/\texttt{opacity} untuk performa optimal.
		\item Custom properties untuk design system dan dark mode.
		\item Aksesibilitas wajib: indikator fokus, reduced motion, kontras warna.
	\end{itemize}
\end{frame}

\begin{frame}{}
	\centering
	\vfill
	\textbf{\Large Pertanyaan?}
	\vfill
\end{frame}

\end{document}
